\section{Related Work}
\paragraph{All-in-one and lightweight restoration.}
TransWeather~\cite{Valanarasu_2022_CVPR} and PromptIR~\cite{NEURIPS2023_e187897e} handle multiple degradations with transformer decoding and input-dependent prompts, respectively. Restormer~\cite{Zamir_2022_CVPR} mixes channels with transposed attention and uses gated feed-forward layers. RAMiT~\cite{Choi_2024_CVPR} combines spatial and channel attention with a small parameter budget. Our comparison measures six trained systems under the same test protocol and records their different initialization sources.

\paragraph{Task utility and precision.}
Task-oriented restoration directly considers the downstream prediction objective~\cite{Chen_2025_CVPR,Gupta_2024_WACV}. We hold the detector fixed to isolate input-view effects. Frequency supervision~\cite{Jiang_2021_FFL} enters a controlled loss ablation. Hardware-aware precision policies~\cite{Wang_2019_HAQ}, integer arithmetic~\cite{Jacob_2018_CVPR}, and transformer post-training quantization~\cite{Yuan_2022_PTQ4ViT,Li_2023_RepQ} address separate parts of the path from numerical sensitivity to executable low-bit inference. Recent quantization work also studies layer sensitivity and restoration-specific fidelity~\cite{Ranjan_2026_MixQViT,Thanh_2026_QATIE,Kim_2026_MPQ}. The NTIRE Efficient Super-Resolution challenge reports runtime alongside parameter count and reconstruction quality~\cite{Ren_2025_CVPRW}. We use parameter count for the model comparison and measure weight sensitivity and export agreement separately.
